% 06-ia-developpement.tex — État de l'art 5/5 : IA générative et
% développement de ressources éducatives (≈ 3 pages).

\chapter{IA générative et développement de ressources éducatives}
\label{chap:ia}

\section{Les assistants de programmation}

\begin{arediger}{1,5 page}
Modèles de langage entraînés sur du code \parencite{chen2021codex} ;
assistants agentiques capables de lire un dépôt, exécuter des tests et
déléguer à des sous-agents. Forces (vitesse, couverture) et limites
(hallucinations, choix implicites, vérification nécessaire).
\end{arediger}

\section{IA et éducation}

\begin{arediger}{1,5 page}
Opportunités et risques pour l'enseignement \parencite{kasneci2023}. Position de
ce travail : l'IA n'est pas utilisée par les étudiants mais par
l'enseignant pour produire des ressources, sous contrôle de règles écrites,
de tests automatiques et de vérifications physiques indépendantes
(section~\ref{sec:methodologie}).
\end{arediger}
